% Options for packages loaded elsewhere
\PassOptionsToPackage{unicode}{hyperref}
\PassOptionsToPackage{hyphens}{url}
%
\documentclass[
  11pt,
  letterpaper,
]{article}
\usepackage{amsmath,amssymb}
\usepackage{iftex}
\ifPDFTeX
  \usepackage[T1]{fontenc}
  \usepackage[utf8]{inputenc}
  \usepackage{textcomp} % provide euro and other symbols
\else % if luatex or xetex
  \usepackage{unicode-math} % this also loads fontspec
  \defaultfontfeatures{Scale=MatchLowercase}
  \defaultfontfeatures[\rmfamily]{Ligatures=TeX,Scale=1}
\fi
\usepackage{lmodern}
\ifPDFTeX\else
  % xetex/luatex font selection
\fi
% Use upquote if available, for straight quotes in verbatim environments
\IfFileExists{upquote.sty}{\usepackage{upquote}}{}
\IfFileExists{microtype.sty}{% use microtype if available
  \usepackage[]{microtype}
  \UseMicrotypeSet[protrusion]{basicmath} % disable protrusion for tt fonts
}{}
\makeatletter
\@ifundefined{KOMAClassName}{% if non-KOMA class
  \IfFileExists{parskip.sty}{%
    \usepackage{parskip}
  }{% else
    \setlength{\parindent}{0pt}
    \setlength{\parskip}{6pt plus 2pt minus 1pt}}
}{% if KOMA class
  \KOMAoptions{parskip=half}}
\makeatother
\usepackage{xcolor}
\usepackage[margin=0.85in]{geometry}
\usepackage{color}
\usepackage{fancyvrb}
\newcommand{\VerbBar}{|}
\newcommand{\VERB}{\Verb[commandchars=\\\{\}]}
\DefineVerbatimEnvironment{Highlighting}{Verbatim}{commandchars=\\\{\}}
% Add ',fontsize=\small' for more characters per line
\newenvironment{Shaded}{}{}
\newcommand{\AlertTok}[1]{\textcolor[rgb]{1.00,0.00,0.00}{\textbf{#1}}}
\newcommand{\AnnotationTok}[1]{\textcolor[rgb]{0.38,0.63,0.69}{\textbf{\textit{#1}}}}
\newcommand{\AttributeTok}[1]{\textcolor[rgb]{0.49,0.56,0.16}{#1}}
\newcommand{\BaseNTok}[1]{\textcolor[rgb]{0.25,0.63,0.44}{#1}}
\newcommand{\BuiltInTok}[1]{\textcolor[rgb]{0.00,0.50,0.00}{#1}}
\newcommand{\CharTok}[1]{\textcolor[rgb]{0.25,0.44,0.63}{#1}}
\newcommand{\CommentTok}[1]{\textcolor[rgb]{0.38,0.63,0.69}{\textit{#1}}}
\newcommand{\CommentVarTok}[1]{\textcolor[rgb]{0.38,0.63,0.69}{\textbf{\textit{#1}}}}
\newcommand{\ConstantTok}[1]{\textcolor[rgb]{0.53,0.00,0.00}{#1}}
\newcommand{\ControlFlowTok}[1]{\textcolor[rgb]{0.00,0.44,0.13}{\textbf{#1}}}
\newcommand{\DataTypeTok}[1]{\textcolor[rgb]{0.56,0.13,0.00}{#1}}
\newcommand{\DecValTok}[1]{\textcolor[rgb]{0.25,0.63,0.44}{#1}}
\newcommand{\DocumentationTok}[1]{\textcolor[rgb]{0.73,0.13,0.13}{\textit{#1}}}
\newcommand{\ErrorTok}[1]{\textcolor[rgb]{1.00,0.00,0.00}{\textbf{#1}}}
\newcommand{\ExtensionTok}[1]{#1}
\newcommand{\FloatTok}[1]{\textcolor[rgb]{0.25,0.63,0.44}{#1}}
\newcommand{\FunctionTok}[1]{\textcolor[rgb]{0.02,0.16,0.49}{#1}}
\newcommand{\ImportTok}[1]{\textcolor[rgb]{0.00,0.50,0.00}{\textbf{#1}}}
\newcommand{\InformationTok}[1]{\textcolor[rgb]{0.38,0.63,0.69}{\textbf{\textit{#1}}}}
\newcommand{\KeywordTok}[1]{\textcolor[rgb]{0.00,0.44,0.13}{\textbf{#1}}}
\newcommand{\NormalTok}[1]{#1}
\newcommand{\OperatorTok}[1]{\textcolor[rgb]{0.40,0.40,0.40}{#1}}
\newcommand{\OtherTok}[1]{\textcolor[rgb]{0.00,0.44,0.13}{#1}}
\newcommand{\PreprocessorTok}[1]{\textcolor[rgb]{0.74,0.48,0.00}{#1}}
\newcommand{\RegionMarkerTok}[1]{#1}
\newcommand{\SpecialCharTok}[1]{\textcolor[rgb]{0.25,0.44,0.63}{#1}}
\newcommand{\SpecialStringTok}[1]{\textcolor[rgb]{0.73,0.40,0.53}{#1}}
\newcommand{\StringTok}[1]{\textcolor[rgb]{0.25,0.44,0.63}{#1}}
\newcommand{\VariableTok}[1]{\textcolor[rgb]{0.10,0.09,0.49}{#1}}
\newcommand{\VerbatimStringTok}[1]{\textcolor[rgb]{0.25,0.44,0.63}{#1}}
\newcommand{\WarningTok}[1]{\textcolor[rgb]{0.38,0.63,0.69}{\textbf{\textit{#1}}}}
\setlength{\emergencystretch}{3em} % prevent overfull lines
\providecommand{\tightlist}{%
  \setlength{\itemsep}{0pt}\setlength{\parskip}{0pt}}
\setcounter{secnumdepth}{5}
\ifLuaTeX
  \usepackage{selnolig}  % disable illegal ligatures
\fi
\usepackage{bookmark}
\IfFileExists{xurl.sty}{\usepackage{xurl}}{} % add URL line breaks if available
\urlstyle{same}
\hypersetup{
  pdftitle={A rare-event obstruction to the uniform exponential GEPP tail in IE-04},
  hidelinks,
  pdfcreator={LaTeX via pandoc}}

\title{A rare-event obstruction to the uniform exponential GEPP tail in
IE-04}
\usepackage{etoolbox}
\makeatletter
\providecommand{\subtitle}[1]{% add subtitle to \maketitle
  \apptocmd{\@title}{\par {\large #1 \par}}{}{}
}
\makeatother
\subtitle{An explicit robust-growth construction and a Gaussian
small-box bound}
\author{}
\date{September 2026}

\begin{document}
\maketitle

\section{Statement and conclusion}\label{statement-and-conclusion}

IE-04 asks for universal \(c_1,c_2>0\) such that, for every admissible
dimension, deterministic center, noise level, and \(x\geq1\),

\[
\Pr\{\rho_{\rm PP}(\bar A+\sigma G)>x(n/\sigma)^{c_1}\}
\leq 2^{-c_2x},
\tag{1}
\]

where \(\|\bar A\|_2\leq1\), \(0<\sigma\leq1\), and \(G\) has
independent real standard normal entries {[}1{]}. The growth factor is
the largest absolute entry among all active Schur complements, divided
by the largest absolute input entry. Elimination is in exact arithmetic.

\textbf{Theorem.} For every \(n\geq2\) and either \(C=0\) or \(C=I_n\),
a standard Gaussian matrix \(G\) satisfies

\[
\Pr\!\left\{\rho_{\rm PP}(C+G)>
\frac12\left(\frac32\right)^{n-1}\right\}
\geq 2^{-n^2(n^2+n+5)}
\geq 2^{-3n^4}.
\tag{2}
\]

Consequently, (1) is false for every choice of positive constants
\(c_1,c_2\). The case \(\bar A=I_n\), \(\sigma=1\) already contradicts
it. This deterministic center is nonsingular and has spectral norm one.

The point is the requirement of an exponential tail in \textbf{every}
\(x\geq1\). An exponentially large growth factor has at least an
inverse-exponential-in-a-polynomial-of-\(n\) probability; the proposed
right-hand side at the corresponding \(x\) is much smaller. This does
not deny a polynomial growth bound with high probability, or settle a
different tail estimate restricted to a smaller range of \(x\).

\section{A strictly pivoted high-growth
matrix}\label{a-strictly-pivoted-high-growth-matrix}

Define \(W_n=(w_{ij})\) by

\[
w_{ij}=\begin{cases}
1,&j=n,\\
1,&j<n\text{ and }i=j,\\
-1/2,&j<n\text{ and }i>j,\\
0,&j<n\text{ and }i<j.
\end{cases}
\tag{3}
\]

For example, its order-four instance is

\[
W_4=\begin{pmatrix}
1&0&0&1\\
-1/2&1&0&1\\
-1/2&-1/2&1&1\\
-1/2&-1/2&-1/2&1
\end{pmatrix}.
\]

At elimination step \(k<n\), the pivot is \(1\) and every competitor in
the pivot column is \(-1/2\). Thus partial pivoting uniquely chooses the
current row. The other nonfinal columns keep the same triangular
pattern. In the final column, all remaining entries have a common value
\(c_k\), and a step changes this to \(c_k-(-1/2)c_k=(3/2)c_k\). Hence

\[
c_k=\left(\frac32\right)^{k-1},\qquad
\rho_{\rm PP}(W_n)=c_n.
\tag{4}
\]

In particular, \(W_n\) is nonsingular. All entries of every exact active
matrix have absolute value at most \(2^n\).

\section{Quantitative robustness on a full entrywise
box}\label{quantitative-robustness-on-a-full-entrywise-box}

Set

\[
B_n=2^{n+2},\qquad \delta_n=2^{-(n^2+n+1)}.
\tag{5}
\]

These constants satisfy

\[
B_n^{n-1}\delta_n=\frac18.
\tag{6}
\]

\textbf{Lemma.} Every real matrix \(A\) with
\(\|A-W_n\|_{\max}\leq\delta_n\) is nonsingular, has a unique no-swap
partial-pivoting path, and satisfies

\[
\rho_{\rm PP}(A)>\frac{c_n}{2}.
\tag{7}
\]

\textbf{Proof.} Denote the active Schur complements of \(W_n\) and \(A\)
by \(S_k\) and \(T_k\). We prove by induction that the same no-swap path
is selected and

\[
e_k:=\|T_k-S_k\|_{\max}\leq B_n^{k-1}\delta_n\leq\frac18.
\tag{8}
\]

The initial claim follows from the box condition. Suppose it holds at
step \(k<n\); retain the original global indices on the active entries.
The perturbed pivot has value at least \(1-e_k\geq7/8\), while every
competing entry has absolute value at most \(1/2+e_k\leq5/8\). The pivot
is therefore strictly largest and positive.

Write \(b=-1/2\) for the unperturbed elimination multiplier in a
remaining row, and \(q=(T_k)_{ik}/(T_k)_{kk}\) for its perturbed value.
For \(e=e_k\leq1/8\),

\[
|q|\leq\frac{1/2+e}{1-e}\leq\frac57<1,
\qquad
|q-b|\leq\frac{(3/2)e}{1-e}\leq2e.
\tag{9}
\]

For each remaining \((i,j)\), subtracting the two exact update formulas
gives

\[
\begin{aligned}
(T_{k+1}-S_{k+1})_{ij}
={}&(T_k-S_k)_{ij}
-q(T_k-S_k)_{kj}\\
&-(q-b)(S_k)_{kj}.
\end{aligned}
\tag{10}
\]

Using \(|(S_k)_{kj}|\leq2^n\) and (9),

\[
e_{k+1}\leq (2+2^{n+1})e_k\leq B_ne_k.
\tag{11}
\]

This proves the induction, including the strict pivot selection at every
nonfinal step. The final scalar Schur complement is at least
\(c_n-1/8>0\), so all pivots are nonzero and \(A\) is nonsingular. Also
\(\|A\|_{\max}\leq1+\delta_n\leq9/8\). Consequently,

\[
\rho_{\rm PP}(A)\geq\frac{c_n-1/8}{1+\delta_n}
\geq\frac{8c_n-1}{9}\geq\frac{7c_n}{9}>\frac{c_n}{2},
\tag{12}
\]

where \(c_n\geq1\) was used. This proves the lemma. Notice that the
statement concerns the entire box, not sampled perturbations, and uses
no unspecified continuity radius.

\section{Gaussian probability of the
box}\label{gaussian-probability-of-the-box}

Fix \(C\in\{0,I_n\}\) and set \(A=C+G\). Define the event

\[
E_n=\{\,|G_{ij}-(w_{ij}-C_{ij})|\leq\delta_n
\text{ for every }i,j\,\}.
\tag{13}
\]

On \(E_n\), the matrix \(A=C+G\) belongs to the box around \(W_n\) in
the lemma. For the smoothed counterexample we use \(C=I_n\), so
\(\bar A=I_n\) and \(\sigma=1\). This obeys \(\|\bar A\|_2=1\) and even
a requirement that the deterministic center be nonsingular. The matrix
\(W_n\) is only the center of the high-growth event in input space;
there is no requirement that \(\|W_n\|_2\leq1\). The case \(C=0\) gives
an additional pure-Gaussian statement but is not needed for the
counterexample.

Each interval in (13) lies within \([-2,2]\), because
\(|w_{ij}-C_{ij}|\leq1\) for both choices of \(C\), and
\(\delta_n\leq1/8\). The standard normal density there is bounded below
by

\[
\frac{e^{-2}}{\sqrt{2\pi}}>\frac1{32}.
\tag{14}
\]

For instance, \(e<3\) and \(\pi<4\) imply that the left side exceeds
\(1/27\). Each interval has length \(2\delta_n\), and the entries are
independent. It follows that

\[
\Pr(E_n)\geq\left(\frac{\delta_n}{16}\right)^{n^2}
=2^{-n^2(n^2+n+5)}.
\tag{15}
\]

The lemma gives \(E_n\subseteq\{\rho_{\rm PP}(C+G)>c_n/2\}\). Finally,
\(n^2(n^2+n+5)\leq3n^4\) for \(n\geq2\): subtracting gives
\(n^2(2n^2-n-5)>0\). This proves (2).

\section{Contradiction for arbitrary proposed
constants}\label{contradiction-for-arbitrary-proposed-constants}

Let \(c_1,c_2>0\) be arbitrary. Put

\[
x_n=\frac{(3/2)^{n-1}}{2n^{c_1}},
\qquad K_n=n^2(n^2+n+5).
\tag{16}
\]

Exponential growth dominates every fixed power of \(n\), so \(x_n\geq1\)
and \(c_2x_n>K_n\) for all sufficiently large \(n\). More explicitly,

\[
\log\frac{x_n}{K_n}
=(n-1)\log(3/2)-(c_1+4)\log n-\log2
-\log(1+1/n+5/n^2)\longrightarrow+\infty.
\tag{17}
\]

For such \(n\), (2) with \(C=I_n\) yields

\[
\Pr\{\rho_{\rm PP}(I_n+G)>x_nn^{c_1}\}
\geq2^{-K_n}>2^{-c_2x_n},
\tag{18}
\]

contradicting (1) in an admissible instance. Thus no pair of universal
positive constants can satisfy the displayed IE-04 inequality.

\section{Verification, scope, and
provenance}\label{verification-scope-and-provenance}

This is an analytic proof for every dimension and every proposed pair of
constants. The accompanying \texttt{verify\_bounds.py} supplies
additional exact checks: the default run certifies the \textbf{whole}
entrywise box for each \(n=2,\ldots,16\) by outward-rounded dyadic
interval arithmetic, verifies the scalar bounds for \(n=2,\ldots,256\),
and gives illustrative exact contradiction dimensions for several
selected constants. These finite checks do not replace the universal
argument above. They require only Python's standard library.

\begin{Shaded}
\begin{Highlighting}[]
\NormalTok{python verify\_bounds.py}
\end{Highlighting}
\end{Shaded}

\texttt{certificate.json} and \texttt{verification.txt} record the
successful run. The earlier interrupted work contained the
high-growth/small-box proof outline; this explicit version, including
the constants and all inequalities, was completed and checked during
recovery. The argument is submitted as an AI-assisted proof for review.
No external human peer review, formal proof-assistant certification, or
novelty/priority determination is claimed.

The conclusion concerns precisely (1), with its unrestricted \(x\geq1\).
It does not produce an alternative optimal smoothed tail, resolve growth
expectations, or identify a sharp high-probability polynomial exponent.
The historical motivation appears in {[}2{]}, Conjecture 16; the
canonical repository formulation is the target certified here.

\section{References}\label{references}

{[}1{]} OpenProblemsInNLA, IE-04, canonical statement, displayed
smoothed probability inequality and growth convention; accessed during
the recovery session.
\href{https://github.com/ajt60gaibb/OpenProblemsInNLA/tree/main/linear-systems-and-elimination/IE-04}{Repository
entry}.

{[}2{]} D. A. Spielman and S.-H. Teng, \emph{Smoothed Analysis of
Algorithms and Heuristics: Progress and Open Questions}, in
\emph{Foundations of Computational Mathematics, Santander 2005} (2006),
274--342. Section P6, Conjecture 16, page 52 of the
\href{https://www.cs.yale.edu/homes/spielman/PAPERS/focmSmoothed.pdf}{author
PDF}. DOI: 10.1017/CBO9780511721571.010.

\end{document}
